\section[Alternative Riemann-Hilbert problem for Painlev\'e-III(D\_8)]{Alternative Riemann--Hilbert problem for Painlev\'e-III($\boldsymbol{D_8}$)}

\subsection[Fabry-type transformation and existence of Rhateven/odd(lambda,z)]{Fabry-type transformation and existence of $\boldsymbol{\widehat{R}^\mathrm{even/odd}(\lambda,z)}$} \label{sec:suleimanov-solution-connection}
The Lax pair \eqref{eq:Lax-system} is unusual in that its coefficient matrices have non-diagonalizable leading terms at both of its singular points $\lambda = 0$ and $\lambda = \infty$, i.e., the coefficients of $\lambda^0$ and $\lambda^{-2}$ in~\eqref{eq:Lambda-exact} are not diagonalizable. To deduce the existence of the matrix functions ${\mathbf{\Omega}}^{\mathrm{even/odd}}(\lambda, z)$ and \smash{$\widehat{\mathbf{R}}^{\mathrm{even/odd}}(\lambda, z)$}, we identify this Lax pair with ones appearing in the literature by considering the following Fabry-type transformation
\begin{align}
&\mathbf{S}(\xi, z) := {\dfrac{1}{\sqrt{2}}} \begin{bmatrix} -\ii & 1\\ -1 &\ii \end{bmatrix} (2z)^{-\sigma_3/4} \xi^{-\sigma_3/2}\nonumber
\\
&\hphantom{\mathbf{S}(\xi, z) :=}{}
\times \begin{cases}
\mathbf \Omega\big(\xi^2 \ee^{\frac{\ii\pi}{2}}, z\big),& -\frac{\pi}{2}<\arg(\xi)<\frac{\pi}{2},\\
\mathbf \Omega\big(\xi^2 \ee^{-\frac{3\pi \ii}{2}}, z\big)(-\ii\sigma_2),& \frac{\pi}{2}<\arg(\xi)<\pi, \\
\mathbf \Omega\big(\xi^2 \ee^{\frac{5\pi \ii}{2}}, z\big)\ii\sigma_2,& -\pi<\arg(\xi)<-\frac{\pi}{2},
 \end{cases}
 \label{eq:fabry-1}
\end{align}
when $|\xi|>1$, and
\begin{gather}
\mathbf{S}(\xi, z) := {\dfrac{1}{\sqrt{2}}} \begin{bmatrix} -\ii & 1\\ -1 &\ii \end{bmatrix} (2z)^{-\sigma_3/4} \xi^{-\sigma_3/2}
\cdot \begin{cases}
\mathbf \Omega\big(\xi^2 \ee^{\frac{\ii\pi}{2}}, z\big),& -\frac{\pi}{2}<\arg(\xi)<\frac{\pi}{2},\\
\mathbf \Omega\big(\xi^2 \ee^{-\frac{3\pi \ii}{2}}, z\big)\ii\sigma_2,& \frac{\pi}{2}<\arg(\xi)<\pi, \\
\mathbf \Omega\big(\xi^2 \ee^{\frac{5\pi \ii}{2}}, z\big)(-\ii\sigma_2),& -\pi<\arg(\xi)<-\frac{\pi}{2},
 \end{cases}\!\!\!\!
 \label{eq:fabry-2}
\end{gather}
when $|\xi|<1$. Denoting
\begin{equation}
\mathbf{K} := {\dfrac{1}{\sqrt{2}}} \begin{bmatrix} \ii & -1\\ 1 & -\ii \end{bmatrix}
\label{eq:K-def}
\end{equation}
and using expansions \eqref{eq:Omega-expand-infty}, \eqref{eq:Omega-expand-zero} (note the branch choices in \eqref{eq:rho-infty-branch}, \eqref{eq:rho-zero-branch}) one can directly check that
\begin{align*}
 \mathbf{S}(\xi, z) = \mathbf{K}^{-1} \Bigg( \mathbb{I} \!+\! \begin{bmatrix} 0 & 0 \\ {(2z)^{1/2}}\Xi^{(8)}_{21}(z) & 0\end{bmatrix} \dfrac{1}{\ii \xi}\! + \!\begin{bmatrix} \Xi^{(8)}_{11}(z) & 0 \\ 0 & {\Xi^{(8)}_{22}(z)} \end{bmatrix} \dfrac{1}{\ii \xi^2}\!+\! \mathcal{O}\big(\xi^{-3}\big)\Bigg)\mathbf{K}\ee^{\ii (2z)^{1/2} \xi \sigma_3},
\end{align*}
as $\xi \to \infty$, and
\begin{align*}
& \mathbf{S}(\xi, z) = \mathbf{K}^{-1} \left( \begin{bmatrix} \Delta^{(8)}_{11}(z) & 0 \\ 0 & 0\end{bmatrix} \dfrac{1}{\xi} + \begin{bmatrix} 0 & \dfrac{\Delta^{(8)}_{12}(z)}{(2z)^{1/2}} \\ (2z)^{1/2} \Delta^{(8)}_{21}(z) & 0 \end{bmatrix} \right.\\
 &\left.\hphantom{\mathbf{S}(\xi, z) =}{}
 + \begin{bmatrix} f(z)& 0 \\ 0 & {\Delta^{(8)}_{22}(z)} \end{bmatrix}\xi + \mathcal{O}\big(\xi^2\big)\right)\ee^{-\frac{\ii\pi}{4}\sigma_3}\mathbf{K} \ee^{ (2z)^{1/2}{ \xi^{-1}} \sigma_3},
\end{align*}
as $\xi\to 0$, where $f(z):= \ii\left((\Delta_{11} \Pi_{11})(z) + (\Delta_{12}\Pi_{21})(z) \right)$; this partly works due to the identity
\begin{equation}
\mathbf{K}\sigma_2 = - \sigma_3 \mathbf K.
\label{eq:K-symmetry}
\end{equation}
Furthermore, one can directly verify that the jump relations \eqref{eq:Omega-jump-circle}--\eqref{eq:Omega-jump-outside-circle} translate to the jumps shown in Figure \ref{fig:D8-limit}, where $\mathbf{C}_{0 \infty}$ is as in \eqref{eq:limiting-connection-matrix}. \begin{figure}
 \centering
 \includegraphics{Figures/6.pdf}
 \caption{The jump contour and matrices for $\widetilde{\mathbf{S}}$, where $\mathbf{S}_1^\infty$, $\mathbf{S}_0^0$ are as in \eqref{eq:Stokes-zero} and $\mathbf{C}_{0\infty}$ is in~\eqref{eq:limiting-connection-matrix}.}
 \label{fig:D8-limit}
\end{figure}
The jump matrices satisfy two cyclic relations about the nonsingular self-intersection points of the jump contour, namely,
 \begin{alignat*}{3}
 &\text{about}\ \xi = +\ii\colon\ &&\mathbf{C}_{0\infty}^{-1} \mathbf{S}_1^\infty (\ii \sigma_2) \mathbf{C}_{0\infty} (\ii \sigma_2) \big(\mathbf{S}_0^0\big)^{-1} = \mathbb{I}, &\\
 &\text{about}\ \xi = -\ii\colon\ &&\left[(\ii \sigma_2) \mathbf{C}_{0\infty} (\ii \sigma_2)\right]^{-1} \big(\mathbf{S}_{0}^0\big)^{-1}\mathbf{C}_{0\infty}\mathbf S_1^\infty = \mathbb{I}.&
\end{alignat*}
Observe that the matrix ${\mathbf{S}}(\xi, z)$ possesses the following useful symmetry:
\begin{equation}
{\mathbf{S}}(\xi, z) = \sigma_2 \begin{cases}
 -{\mathbf{S}}(-\xi, z) \sigma_2, & |\xi|<1, \\
{\mathbf{S}}(-\xi, z) \sigma_2, & |\xi|>1.
\label{eq:S-symmetry}
\end{cases}
\end{equation}
This result also uses the identity \eqref{eq:K-symmetry}. Using this symmetry, it can be checked that the Fabry transformation \eqref{eq:fabry-1}--\eqref{eq:fabry-2} is invertible with{\samepage
\begin{equation}
\mathbf \Omega(\lambda ) = \rho_\infty^{\sigma_3/2}(\lambda, z) \mathbf{K} \mathbf S\big(\sqrt{-\ii \lambda}\big),
\label{eq:inverse-fabry}
\end{equation}
where all roots are principal branches.}

While the singular behavior of $\mathbf{S}(\xi, z)$ at $\xi = 0$ is concerning, the fact that the leading coefficient is a singular matrix allows us to handle this problem by letting
\begin{equation}\label{eq:triangular-transformation}
\widetilde{\mathbf{S}}(\xi, z) = \left( \mathbb{I} -\dfrac{1}{\xi} \mathbf{T}(z) \right) \mathbf{S}(\xi, z),
\end{equation}
where
\begin{equation}
\mathbf{T}(z) = \mathbf{K}^{-1}\begin{bmatrix} 0 & \dfrac{\Delta^{(8)}_{11}(z)}{(2z)^{1/2}\Delta^{(8)}_{21}(z)} \\ 0 & 0\end{bmatrix}\mathbf{K}.
\label{eq:L-def}
\end{equation}
Since the prefactor is analytic in $\C \setminus \{0\}$, the jumps of $\widetilde{\mathbf{S}}$ are unchanged. As for the asymptotic behavior, as $\xi\to\infty$
\begin{align*}
\widetilde{\mathbf S}(\xi, z) = \left( \mathbb{I} + \mathbf K^{-1} \begin{bmatrix}
 0 & -\dfrac{\Delta^{(8)}_{11}(z)}{(2z)^{1/2}\Delta^{(8)}_{21}(z)} \\ -\ii(2z)^{1/2} \Xi^{(8)}_{21}(z) & 0
\end{bmatrix} \mathbf K \dfrac{1}{\xi}+ \mathcal{O}\big(\xi^{-2}\big) \right)\ee^{\ii (2z)^{1/2} \xi \sigma_3},
\end{align*}
and as $\xi\to 0$
\begin{align*}
\widetilde{\mathbf S}(\xi, z) = \left( \mathbf{K}^{-1}\begin{bmatrix} 0 & -\dfrac{1}{(2z)^{1/2}\Delta^{(8)}_{21}(z)} \\ (2z)^{1/2}\Delta^{(8)}_{21}(z) & 0\end{bmatrix} \ee^{-\frac{\ii\pi}{4}\sigma_3}\mathbf{K} + \mathcal{O}(\xi) \right)\ee^{{ (2z)^{1/2}}{ \xi^{-1}} \sigma_3}.
\end{align*}
\begin{Remark}
Noting that
\[
\det \left( \mathbf{K}^{-1}\begin{bmatrix} 0 & -\dfrac{1}{(2z)^{1/2}\Delta^{(8)}_{21}(z)} \\ (2z)^{1/2}\Delta^{(8)}_{21}(z) & 0\end{bmatrix} \ee^{-\frac{\ii\pi}{4}\sigma_3}\mathbf{K} \right) = 1,
\]
one can carry out a computation similar to the one in Section \ref{sec:Lax-pair} to arrive at a pair of differential equations analogous to \eqref{eq:Lax-system}, but with diagonalizable leading matrices at the two singular points at $\xi = 0, \infty$; this system appears in \cite[Chapter 2]{N} and \cite{GL}, for example. Since we do not make use of this Lax pair, we omit the calculation.
\end{Remark}

Using \eqref{eq:K-symmetry}, it follows that
\[
\mathbb I - \dfrac{1}{\xi}\mathbf T(z) = \sigma_2\left( \mathbb I + \dfrac{1}{\xi}\mathbf T(z) \right) \sigma_2,
\]
which implies that matrix $\widetilde{\mathbf{S}}(\xi, z)$ also satisfies the symmetry \eqref{eq:S-symmetry}. To simplify this symmetry, let
\begin{equation}
\widehat{\mathbf{S}}(\xi, z) := \ee^{-\pi \ii \sigma_3/4 } \begin{cases} \widetilde{\mathbf{S}}(\xi, z) \ee^{\pi \ii \sigma_3/4}, & |\xi| >1, \\
\widetilde{\mathbf{S}}(\xi, z) \ee^{-\pi \ii \sigma_3/4}, & |\xi| <1.
\end{cases}
 \label{eq:rotation-transformation}
\end{equation}
Then, $\widehat{\mathbf{S}}(\xi, z)$ solves the following Riemann--Hilbert problem.
\begin{rhp}
Let $(y_1,y_2, y_3) \in \C^3$ be the monodromy data corresponding to~$U(z)$ given in \eqref{eq:y1}--\eqref{eq:y3}, and fix $z\in\mathbb{C}$. Seek a $2\times 2$ matrix function~${\xi\mapsto\widehat{\mathbf{S}}(\xi, z)}$ satisfying the following properties:
\begin{itemize}\itemsep=0pt
 \item Analyticity: $\widehat{\mathbf{S}}(\xi, z)$ is an analytic function of $\xi$ for $|\xi|\neq 1$.
 \item Jump condition:
 $\widehat{\mathbf{S}}(\xi, z)$ takes analytic boundary values on the unit circle from the interior and exterior, denoted $\widehat{\mathbf{S}}_+(\xi,z)$ and $\widehat{\mathbf{S}}_-(\xi,z)$ for $|\xi|=1$ respectively, and they are related by
 \begin{equation*}
 % \label{S-hat-jump}
 \widehat{\mathbf{S}}_+(\xi,z)=\widehat{\mathbf{S}}_-(\xi,z) \mathbf J_{\widehat{\mathbf{S}}}(\xi),
 \end{equation*}
 where $\mathbf J_{\widehat{\mathbf{S}}}(\xi)$ is shown in Figure {\rm\ref{fig:D8-limit-rotated}} and
 \begin{align}
 &\widehat{\mathbf C}_{0\infty} := \ee^{\pi \ii \sigma_3/4} \mathbf C_{0\infty}\ee^{\pi \ii \sigma_3/4}, \qquad
 \widehat{\mathbf S}_1^\infty := \ee^{-\pi \ii \sigma_3/4} {\mathbf S}_1^\infty \ee^{\pi \ii \sigma_3/4} = \ee^{\pi \ii \sigma_3/4} ({\mathbf S}_1^\infty)^{-1} \ee^{-\pi \ii \sigma_3/4}, \nonumber\\
 & \widehat{\mathbf{S}}_0^0:= \ee^{\pi \ii \sigma_3/4} {\mathbf S}_0^0 \ee^{-\pi \ii \sigma_3/4} = \ee^{-\pi \ii \sigma_3/4} \big({\mathbf S}_0^0\big)^{-1} \ee^{\pi \ii \sigma_3/4}. \label{eq:hat-jumps}
 \end{align}
 \item Normalization:
\begin{gather}\label{eq:s-hat-norm-inf}\widehat{\mathbf{S}}(\xi,z) = \big(\mathbb{I} +\widehat{\mathbf{\Xi}}^{(8)}(z)\xi^{-1}+ \mathcal{O}\big(\xi^{-2}\big) \big) \ee^{\ii (2z)^{1/2} \xi \sigma_3 } \qquad \text{as} \quad \xi \to \infty,
\end{gather}
 and
\begin{gather}
\widehat{\mathbf{S}}(\xi,z) = \widehat{\mathbf{\Delta}}^{(8)}(z) \big(\mathbb{I} + \widehat{\mathbf{\Pi}}(z)\xi+\mathcal{O}\big(\xi^2\big) \big) \ee^{ (2z)^{1/2} \xi^{-1} \sigma_3 } \qquad \text{as} \quad \xi \to 0,\label{eq:s-hat-norm-zero}
\end{gather}
 where $\widehat{\mathbf{\Delta}}^{(8)}(z)$ may be written in terms of entries of $\mathbf{\Delta}^{(8)}(z)$ and $\mathbf{K}$.
\end{itemize}
\label{rhp:S-hat}
\end{rhp}
\begin{figure}
 \centering
 \includegraphics{Figures/7.pdf}
 \caption{The jump contour and matrices for $\hat{\mathbf S}$, where the jump matrices are as in \eqref{eq:hat-jumps}.}
 \label{fig:D8-limit-rotated}
\end{figure}
Now, the matrix $\widehat{\mathbf{S}}(\xi,z)$ satisfies the symmetry
\[
\sigma_1 \widehat{\mathbf{S}}(-\xi,z) \sigma_1 = \widehat{\mathbf{S}}(\xi,z) .
\]
Furthermore, it was shown in \cite[Theorem 4]{N} that matrix $\widehat{\mathbf{S}}(\xi, z)$ exists for all $z$ outside of a~discrete set $\Sigma$ and is a meromorphic function of $z$ in $\C \setminus \Sigma$. Since the transformations used to arrive to $\widehat{\mathbf{S}}(\xi,z) $ from $\widehat{\mathbf{R}}(\lambda,z) $ and ${\mathbf{\Omega}}(\lambda,z) $ are invertible, we deduce the existence of matrix functions satisfying Riemann--Hilbert Problems \ref{rhp:Rhat-even/odd} and \ref{rhp:D8}.

It was shown in \cite{Palmer} that $\Sigma$ coincides with the set of zeros of the $\tau$-function associated to the Riemann--Hilbert problem. According to \cite{JMU} the expression for the logarithmic derivative of the $\tau$-function associated to Riemann--Hilbert Problem \ref{rhp:S-hat} is given by
\begin{gather} \label{eq:jmu-tau}
\dfrac{\mathrm{d}}{\mathrm{d}z}\ln(\tau(z))=-\frac{1}{\sqrt{2z}}\big(\Tr\big(\widehat{\mathbf{\Pi}}(z)\sigma_3\big)+\ii\Tr\big(\widehat{\mathbf{\Xi}}^{(8)}(z)\sigma_3\big)\big).
\end{gather}
 After a long symbolic computation using the transformations \eqref{eq:fabry-1}, \eqref{eq:fabry-2}, \eqref{eq:triangular-transformation}, and \eqref{eq:rotation-transformation}, we get the following result. The Lax pair \eqref{eq:lax_pair_D8}--\eqref{eq:lax_pair_D8_2} is gauge equivalent to
\[
\frac{\partial\widehat{\mathbf{S}}}{\partial\xi} (\xi, z) = {\widehat{\mathbf{\Lambda}}}^{(8)}(\xi, z) \widehat{\mathbf{S}}(\xi, z), \qquad
\frac{\partial\widehat{\mathbf{S}}}{\partial z} (\xi, z) = {\widehat{\mathbf{Z}}}(\xi, z) \widehat{\mathbf{S}}(\xi, z),
\]
where
\begin{align*}
&\widehat{\mathbf{\Lambda}}^{(8)}(\xi, z) = \ii(2z)^{1/2}\sigma_3 +\frac{1}{\xi}\frac{zU'(z)}{2U(z)}\sigma_1+\frac{\ii}{\xi^2}\left({\frac{z}{2}}\right)^{1/2}\left(U(z)-\frac{1}{U(z)}\right)\sigma_3\\
&\hphantom{\widehat{\mathbf{\Lambda}}^{(8)}(\xi, z) =}{}
+\frac{1}{\xi^2}\left({\frac{z}{2}}\right)^{1/2}\left(U(z)+\frac{1}{U(z)}\right)\sigma_2,
\end{align*}
and
\begin{align*}
&\widehat{\mathbf{Z}}(\xi, z) = \frac{\ii\xi}{(2z)^{1/2}}\sigma_3 +\frac{U'(z)}{4U(z)}\sigma_1-\frac{\ii}{2\xi}\frac{1}{(2z)^{1/2}}\left(U(z)-\frac{1}{U(z)}\right)\sigma_3\\
&\hphantom{\widehat{\mathbf{Z}}(\xi, z) =}{}
-\frac{1}{2\xi}\frac{1}{(2z)^{1/2}}\left(U(z)+\frac{1}{U(z)}\right)\sigma_2.
\end{align*}
Similarly, the coefficients in \eqref{eq:s-hat-norm-inf}--\eqref{eq:s-hat-norm-zero} have the following expressions
\begin{gather}
\widehat{\mathbf{\Delta}}^{(8)}(z)=\big(\ee^{-\ii\pi/4}U(z)^{1/2}\big)^{\sigma_1},\nonumber\\
\widehat{\mathbf{\Xi}}^{(8)}(z)=\ii\left(\frac{z}{2}\right)^{1/2}\left(\frac{zU'(z)^2}{8U(z)^2}-
 U(z)+\frac1{U(z)}\right)\sigma_3-\left(\frac{z}{2}\right)^{1/2}\frac{U'(z)}{4U(z)}\sigma_2,\nonumber\\
 \widehat{\mathbf{\Pi}}(z)=-\left(\frac{z}{2}\right)^{1/2}\left(\frac{zU'(z)^2}{8U(z)^2}-
 U(z)+\frac1{U(z)}\right)\sigma_3+\ii\left(\frac{z}{2}\right)^{1/2}\frac{U'(z)}{4U(z)}\sigma_2.\label{eq:s-hat-residues}
\end{gather}
In our computation we expressed $W(z)$, $X(z)$, and $V(z)$ in terms of $U(z)$ and $U'(z)$ using the identities \eqref{eq:UVWXidentity1}--\eqref{eq:UVWidentity1} and the first equation in \eqref{eq:PossiblyD8system}. Plugging \eqref{eq:s-hat-residues} into \eqref{eq:jmu-tau} we get
\begin{gather*}
\dfrac{\mathrm{d}}{\mathrm{d}z}\ln(\tau(z))=\frac{zU'(z)^2}{4U(z)^2}-
 2 U(z)+\frac2{U(z)}.
\end{gather*}
Differentiating once again, we have
\begin{gather} \label{eq:d8-solvability}
\dfrac{\mathrm{d}^2}{\mathrm{d}z^2}\ln(\tau(z))=-\frac{1}{4}\left(\dfrac{\mathrm{d}}{\mathrm{d}z}\ln(U(z))\right)^2.
\end{gather}
 Now we see from \eqref{eq:d8-solvability} that the set $\Sigma$ of zeros of the $\tau$-function coincides precisely with the union of poles and zeros of the function $U(z)$.

\subsection[Relationship between U\^{}\{even\}, U\^{}\{odd\}]{Relationship between $\boldsymbol{U^{\mathrm{even}}}$, $\boldsymbol{U^{\mathrm{odd}}}$}

To complete the proof of Theorem \ref{thm:general}, we must show that $U^{\mathrm{odd}}(z) = -1/U^{\mathrm{even}}(z)$. One can already observe that this should be the case by checking that the leading behavior predicted in Theorem \ref{thm:D8-asymptotics-zero} satisfies the involution, but we now present a proof on the level of Riemann--Hilbert problems. First, note that if one chooses the square root in \eqref{eq:Vm-even} in such a way that~$\mathbf{V}^\mathrm{even} = \ii \sigma_3 \mathbf{V}^\mathrm{odd}$, it follows from \eqref{eq:limiting-connection-matrix} that\footnote{One can check that making the other choice of the square root yields the same connection matrix but with the opposite sign, and so it follows from Remark \ref{rem:minus-y1-y2} that this choice is immaterial.}
\[
 \mathbf{C}_{0\infty}^{\mathrm{odd}} = \sigma_3 \mathbf{C}^{\mathrm{even}}_{0 \infty} \sigma_3.
\]
This, in particular, implies the symmetry
\[
\widetilde{\mathbf{S}}^{\mathrm{odd}}(\lambda, z) = \sigma_3 \widetilde{\mathbf{S}}^{\mathrm{even}}(\lambda, z) \sigma_3,
\]
and, in view of \eqref{eq:inverse-fabry}, we have
\begin{equation}
 \mathbf{\Omega}^{\mathrm{odd}}(\lambda, z) = \rho_\infty^{\sigma_3/2}(\lambda, z) \mathbf{K}\left( \mathbb{I} + \dfrac{1}{\sqrt{-\ii \lambda}} \mathbf{T}^{\mathrm{even}}(z) \right) \sigma_3 \widetilde{\mathbf{S}}^{\mathrm{even}}\big(\sqrt{-\ii \lambda}, z\big) \sigma_3.
 \label{eq:Omega-even-odd}
\end{equation}
Recalling \eqref{eq:rho-infty-branch}, \eqref{eq:K-def}, \eqref{eq:L-def}, and the identity $2\ii U(z)X(z) = \Delta^{(8)}_{11}(z)/\Delta^{(8)}_{21}(z)$, see \eqref{eq:Delta-U-X-identity}, we have~$\mathbf{\Omega}^{\mathrm{odd}}(\lambda, z) = \mathbf{G}(\lambda, z) \mathbf{\Omega}^{\mathrm{even}}(\lambda, z)$, where
\begin{align}
\mathbf{G}(\lambda, z):={}& \rho_\infty^{\sigma_3/2}(\lambda, z) \mathbf{K}\left( \mathbb{I} + \dfrac{1}{\sqrt{-\ii \lambda}} \mathbf{T}^{\mathrm{even}}(z) \right) \sigma_3 \left( \mathbb{I} + \dfrac{1}{\sqrt{-\ii \lambda}} \mathbf{T}^{\mathrm{even}}(z) \right) \rho_\infty^{-\sigma_3/2}(\lambda, z) \nonumber\\
={}& \frac{1}{\rho_\infty(\lambda, z)}\begin{bmatrix} 2U^{\mathrm{even}}(z)X^{\mathrm{even}}(z) & -4\ii (U^{\mathrm{even}}(z) X^{\mathrm{even}}(z))^2 \\ -\ii & -2U^{\mathrm{even}}(z)X^{\mathrm{even}}(z) \end{bmatrix}\nonumber\\
& + \rho_\infty(\lambda, z) \begin{bmatrix}
 0 & \ii \\ 0 & 0
\end{bmatrix}.\label{eq:G-def}
\end{align}
To deduce the relationship between $U^{\mathrm{even}}$, $U^{\mathrm{odd}}$, we now recall that $\mathbf{\Omega}^{\mathrm{even/odd}}$ satisfy the Lax pair~\eqref{eq:Lax-system}. Transforming $\mathbf{\Omega}^{\mathrm{even}}$ as in the right-hand side of \eqref{eq:Omega-even-odd} induces a gauge transformation of the $\lambda$-equation and we have that $\mathbf{\Omega}^{\mathrm{odd}}$ satisfies two equations; the first is the one in \eqref{eq:Lax-system} and the second is
\[
\dpd{\mathbf{\Omega}^{\mathrm{odd}}}{\lambda}(\lambda, z) = \widetilde{\mathbf{\Lambda}}(\lambda, z) \mathbf{\Omega}^{\mathrm{odd}}(\lambda, z),
\]
where
\[
\widetilde{\mathbf{\Lambda}}(\lambda, z) = \dpd{\mathbf{G}}{\lambda}(\lambda, z) \mathbf{G}^{-1}(\lambda, z) + \mathbf{G}(\lambda, z) \mathbf{\Lambda}^{\mathrm{even}}(\lambda, z)\mathbf{G}^{-1}(\lambda, z).
\]
Using \eqref{eq:G-def} and \eqref{eq:UVWXidentity1}, we see that
\begin{align*}
& \widetilde{\mathbf{\Lambda}}(\lambda, z) = \begin{bmatrix} 0 & \ii z \\ 0 & 0 \end{bmatrix}\\
 & \hphantom{\widetilde{\mathbf{\Lambda}}(\lambda, z) =}{}
 + \dfrac{1}{4 \!\lambda}\begin{bmatrix} 2 - V^{\mathrm{even}}(z) + 8\ii U^{\mathrm{even}}(z)X^{\mathrm{even}}(z) & F^\mathrm{even}(z) \\ 2 & -2 + V^{\mathrm{even}}(z) - 8\ii U^{\mathrm{even}}(z)X^{\mathrm{even}}(z) \end{bmatrix} \\
 & \hphantom{\widetilde{\mathbf{\Lambda}}(\lambda, z) =}{} -\dfrac{1}{\lambda^2} \!\begin{bmatrix} (U^{\mathrm{even}}(z))^2 X^{\mathrm{even}}(z) & 2\ii (U^{\mathrm{even}}(z))^3(X^{\mathrm{even}}(z))^2 \\ \ii U^{\mathrm{even}}(z)/2 & - (U^{\mathrm{even}}(z))^2 X^{\mathrm{even}}(z)\end{bmatrix},
\end{align*}
where
\[
F^\mathrm{even}(z) := 4\ii U^{\mathrm{even}}(z) X^{\mathrm{even}}(z) \left(V^{\mathrm{even}}(z) + 6 U^{\mathrm{even}}(z) X^{\mathrm{even}}(z)-2 \ii U^{\mathrm{even}}(z) \right)- \dfrac{4z}{ U^{\mathrm{even}}(z)}.
\]
Since $\det(\mathbf{\Omega}^\mathrm{odd}) = 1$, it follows that $\widetilde{\mathbf{\Lambda}}(\lambda, z) = \mathbf{\Lambda}^{\mathrm{odd}}(\lambda, z)$ and we arrive at identities relating all the potentials $U^{\mathrm{even/odd}}(z)$, $V^{\mathrm{even/odd}}(z)$, $W^{\mathrm{even/odd}}(z)$, $X^{\mathrm{even/odd}}(z)$; comparing the (2,1) entries of the coefficient of $\lambda^{-2}$ yields the desired relation
\[
U^{\mathrm{even/odd}}(z) = -1/U^{\mathrm{odd/even}}(z).
\]


